% Don't like 10pt? Try 11pt or 12pt
\documentclass[10pt]{article}
\RequirePackage[T1]{fontenc}

% % LaTeX will typeset using Computer Modern Roman, which a lot of
% % non-mathematicians and non-engineers won't like. Also, a few PDF
% % viewers may not render CMR very well. Instead, Times New Roman can
% % be used. That's what this package does.
% \usepackage{times}

\usepackage{calc}

% This package helps LaTeX auto-hyphenate hyphenated words if you use
% special hyphens. For example, bio\-/mimicry will properly hyphenate
% ``mimicry'' if necessary.
% \usepackage[shortcuts]{extdash}

% Layout: Puts the section titles on left side of page
\reversemarginpar

%
%         PAPER SIZE, PAGE NUMBER, AND DOCUMENT LAYOUT NOTES:
%
% The next \usepackage line changes the layout for CV style section
% headings as marginal notes. It also sets up the paper size as either
% letter or A4. By default, letter was used. If A4 paper is desired,
% comment out the letterpaper lines and uncomment the a4paper lines.
%
% As you can see, the margin widths and section title widths can be
% easily adjusted.
%
% ALSO: Notice that the includefoot option can be commented OUT in order
% to put the PAGE NUMBER *IN* the bottom margin. This will make the
% effective text area larger.
%
% IF YOU WISH TO REMOVE THE ``of LASTPAGE'' next to each page number,
% see the note about the +LP and -LP lines below. Comment out the +LP
% and uncomment the -LP.
%
% IF YOU WISH TO REMOVE PAGE NUMBERS, be sure that the includefoot line
% is uncommented and ALSO uncomment the \pagestyle{empty} a few lines
% below.
%

%% Use these lines for letter-sized paper
\usepackage[paper=letterpaper,
            %includefoot, % Uncomment to put page number above margin
            marginparwidth=1.2in,     % Length of section titles
            marginparsep=.05in,       % Space between titles and text
            margin=1in,               % 1 inch margins
            includemp]{geometry}

%% Use these lines for A4-sized paper
%\usepackage[paper=a4paper,
%            %includefoot, % Uncomment to put page number above margin
%            marginparwidth=30.5mm,    % Length of section titles
%            marginparsep=1.5mm,       % Space between titles and text
%            margin=25mm,              % 25mm margins
%            includemp]{geometry}

%% More layout: Get rid of indenting throughout entire document
\setlength{\parindent}{0in}

% Provides special list environments and macros to create new ones
\usepackage[shortlabels]{enumitem}

% Simpler bibsections for CV sections
% (thanks to natbib for inspiration)
%
% * For lists of references with hanging indents and no numbers:
%
%   \begin{bibsection}
%       \item ...
%   \end{bibsection}
%
% * For numbered lists of references (with hanging indents):
%
%   \begin{bibenum}
%       \item ...
%   \end{bibenum}
%
%   Note that bibenum numbers continuously throughout. To reset the
%   counter, use
%
%   \restartlist{bibenum}
%
%   at the place where you want the numbering to reset.

\makeatletter
\newlength{\bibhang}
\setlength{\bibhang}{1em}
\newlength{\bibsep}
 {\@listi \global\bibsep\itemsep \global\advance\bibsep by\parsep}
\newlist{bibsection}{itemize}{3}
\setlist[bibsection]{label=,leftmargin=\bibhang,%
        itemindent=-\bibhang,
        itemsep=\bibsep,parsep=\z@,partopsep=0pt,
        topsep=0pt}
\newlist{bibenum}{enumerate}{3}
\setlist[bibenum]{label=[\arabic*],resume,leftmargin={\bibhang+\widthof{[999]}},%
        itemindent=-\bibhang,
        itemsep=\bibsep,parsep=\z@,partopsep=0pt,
        topsep=0pt}
\let\oldendbibenum\endbibenum
\def\endbibenum{\oldendbibenum\vspace{-.6\baselineskip}}
\let\oldendbibsection\endbibsection
\def\endbibsection{\oldendbibsection\vspace{-.6\baselineskip}}
\makeatother

%% Reference the last page in the page number
%
% NOTE: comment the +LP line and uncomment the -LP line to have page
%       numbers without the ``of ##'' last page reference)
%
% NOTE: uncomment the \pagestyle{empty} line to get rid of all page
%       numbers (make sure includefoot is commented out above)
%
\usepackage{fancyhdr,lastpage}
\pagestyle{fancy}
%\pagestyle{empty}      % Uncomment this to get rid of page numbers
\fancyhf{}\renewcommand{\headrulewidth}{0pt}
\fancyfootoffset{\marginparsep+\marginparwidth}
\newlength{\footpageshift}
\setlength{\footpageshift}
          {0.5\textwidth+0.5\marginparsep+0.5\marginparwidth-2in}
\lfoot{\hspace{\footpageshift}%
       \parbox{4in}{\, \hfill %
                    \arabic{page} of \protect\pageref*{LastPage} % +LP
%                    \arabic{page}                               % -LP
                    \hfill \,}}

% Finally, give us PDF bookmarks
\usepackage{amssymb}
\usepackage{xcolor,hyperref}
\definecolor{darkblue}{rgb}{0.0,0.0,0.3}
\hypersetup{colorlinks,breaklinks,
            linkcolor=darkblue,urlcolor=darkblue,
            anchorcolor=darkblue,citecolor=darkblue}

%%%%%%%%%%%%%%%%%%%%%%%% End Document Setup %%%%%%%%%%%%%%%%%%%%%%%%%%%%


%%%%%%%%%%%%%%%%%%%%%%%%%%% Helper Commands %%%%%%%%%%%%%%%%%%%%%%%%%%%%

%%% HEADING AT TOP OF CURRICULUM VITAE

% The title (name) with a horizontal rule under it
% (optional argument typesets an object right-justified across from name
%  as well)
%
% Usage: \makeheading{name}
%        OR
%        \makeheading[right_object]{name}
%
% Place at top of document. It should be the first thing.
% If ``right_object'' is provided in the square-braced optional
% argument, it will be right justified on the same line as ``name'' at
% the top of the CV. For example:
%
%       \makeheading[\emph{Curriculum vitae}]{Your Name}
%
% will put an emphasized ``Curriculum vitae'' at the top of the document
% as a title. Likewise, a picture could be included:
%
%   \makeheading[{\includegraphics[height=1.5in]{my_picture}}]{Your Name}
%
% the picture will be flush right across from the name. For this example
% to work, make sure the extra set of curly braces is included. Also
% makes ure that \usepackage{graphicx} is somewhere in the preamble.
\newcommand{\makeheading}[2][]%
        {\hspace*{-\marginparsep minus \marginparwidth}%
         \begin{minipage}[t]{\textwidth+\marginparwidth+\marginparsep}%
             {\large \bfseries #2 \hfill #1}\\[-0.15\baselineskip]%
                 \rule{\columnwidth}{1pt}%
         \end{minipage}}

%%% SECTION HEADINGS

% The section headings. Flush left in small caps down pseudo-margin.
%
% Usage: \section{section name}
\renewcommand{\section}[1]{\pagebreak[3]%
    \vspace{1.3\baselineskip}%
    \phantomsection\addcontentsline{toc}{section}{#1}%
    \noindent\llap{\scshape\smash{\parbox[t]{\marginparwidth}{\hyphenpenalty=10000\raggedright #1}}}%
    \vspace{-\baselineskip}\par}

%%% LISTS

% This macro alters a list by removing some of the space that follows the list
% (is used by lists below)
\newcommand*\fixendlist[1]{%
    \expandafter\let\csname preFixEndListend#1\expandafter\endcsname\csname end#1\endcsname
    \expandafter\def\csname end#1\endcsname{\csname preFixEndListend#1\endcsname\vspace{-0.6\baselineskip}}}

% These macros help ensure that items in outer-type lists do not get
% separated from the next line by a page break
% (they are used by lists below)
\let\originalItem\item
\newcommand*\fixouterlist[1]{%
    \expandafter\let\csname preFixOuterList#1\expandafter\endcsname\csname #1\endcsname
    \expandafter\def\csname #1\endcsname{\let\oldItem\item\def\item{\pagebreak[2]\oldItem}\csname preFixOuterList#1\endcsname}
    \expandafter\let\csname preFixOuterListend#1\expandafter\endcsname\csname end#1\endcsname
    \expandafter\def\csname end#1\endcsname{\let\item\oldItem\csname preFixOuterListend#1\endcsname}}
\newcommand*\fixinnerlist[1]{%
    \expandafter\let\csname preFixInnerList#1\expandafter\endcsname\csname #1\endcsname
    \expandafter\def\csname #1\endcsname{\let\oldItem\item\let\item\originalItem\csname preFixInnerList#1\endcsname}
    \expandafter\let\csname preFixInnerListend#1\expandafter\endcsname\csname end#1\endcsname
    \expandafter\def\csname end#1\endcsname{\csname preFixInnerListend#1\endcsname\let\item\oldItem}}

% An itemize-style list with lots of space between items
%
% Usage:
%   \begin{outerlist}
%       \item ...    % (or \item[] for no bullet)
%   \end{outerlist}
\newlist{outerlist}{itemize}{3}
    \setlist[outerlist]{label=\enskip\textbullet,leftmargin=*}
    \fixendlist{outerlist}
    \fixouterlist{outerlist}

% An environment IDENTICAL to outerlist that has better pre-list spacing
% when used as the first thing in a \section
%
% Usage:
%   \begin{lonelist}
%       \item ...    % (or \item[] for no bullet)
%   \end{lonelist}
\newlist{lonelist}{itemize}{3}
    \setlist[lonelist]{label=\enskip\textbullet,leftmargin=*,partopsep=0pt,topsep=0pt}
    \fixendlist{lonelist}
    \fixouterlist{lonelist}

% An itemize-style list with little space between items
%
% Usage:
%   \begin{innerlist}
%       \item ...    % (or \item[] for no bullet)
%   \end{innerlist}
\newlist{innerlist}{itemize}{3}
\setlist[innerlist]{label={ },leftmargin=20pt,parsep=0pt,itemsep=0pt,topsep=0pt,partopsep=0pt}
    \fixinnerlist{innerlist}

% An environment IDENTICAL to innerlist that has better pre-list spacing
% when used as the first thing in a \section
%
% Usage:
%   \begin{loneinnerlist}
%       \item ...    % (or \item[] for no bullet)
%   \end{loneinnerlist}
\newlist{loneinnerlist}{itemize}{3}
    \setlist[loneinnerlist]{label=\enskip\textbullet,leftmargin=*,parsep=0pt,itemsep=0pt,topsep=0pt,partopsep=0pt}
    \fixendlist{loneinnerlist}
    \fixinnerlist{loneinnerlist}

%%% EXTRA SPACE

% To add some paragraph space between lines.
% This also tells LaTeX to preferably break a page on one of these gaps
% if there is a needed pagebreak nearby.
\newcommand{\blankline}{\quad\pagebreak[3]}
\newcommand{\halfblankline}{\quad\vspace{-0.5\baselineskip}\pagebreak[3]}

%%% FORMATTING MACROS

% Provides a linked \doi{#1} that links doi:#1 to http://dx.doi.org/#1
\usepackage{doi}
% To change the text before the DOI, adjust this command
%\renewcommand\doitext{doi:}

% Provides a linked \url{#1} that doesn't require escape characters
\usepackage{url}

% You can adjust the style \url{} uses here:
% (options are: same, rm, sf, tt; defaults to tt)
\urlstyle{same}

% For \email{ADDRESS}, links ADDRESS to the url mailto:ADDRESS
% (uncomment to typeset the e\-/mail address in typewriter font;
%  otherwise, will be typeset in the \urlstyle above)
%\DeclareUrlCommand\emaillink{\urlstyle{tt}}
\providecommand*\emaillink[1]{\nolinkurl{#1}}
\providecommand*\email[1]{\href{mailto:#1}{\emaillink{#1}}}

\providecommand\BibTeX{{B\kern-.05em{\sc i\kern-.025em b}\kern-.08em \TeX}}
\providecommand\Matlab{\textsc{Matlab}}

% Custom hyphenation rules for words that LaTeX has trouble with
\hyphenation{bio-mim-ic-ry bio-in-spi-ra-tion re-us-a-ble pro-vid-er Media-Wiki}

%%% RESEARCH INTERESTS BANNER
% A full-width tinted panel (spanning the section-title margin too) that
% breaks from the margin-label layout used by every other section.
% Two interests flank a central two-way arrow.
%
% Usage:
%   \begin{interests}{tagline}
%     \interest{heading}{topic}{description}\interestbridge\interest{...}{...}{...}
%   \end{interests}
\newlength{\rifull}
\newlength{\riinner}
\newlength{\ricol}
\newlength{\ribridge}\setlength{\ribridge}{3em}
\newcommand{\interest}[3]{%
    \begin{minipage}[t]{\ricol}\centering\small
        {\color{darkblue}\scshape #1}\\[1pt]
        {\footnotesize\itshape #2}\\[4pt]
        #3
    \end{minipage}}
\newcommand{\interestbridge}{%
    \makebox[\ribridge]{\raisebox{-0.9\baselineskip}{\color{darkblue}\Large$\rightleftarrows$}}}
\newsavebox{\ribox}
\newenvironment{interests}[1]
    {\par\vspace{1.6\baselineskip}\pagebreak[3]%
     \setlength{\rifull}{\textwidth+\marginparwidth+\marginparsep}%
     \setlength{\fboxsep}{10pt}%
     \setlength{\riinner}{\rifull-2\fboxsep}%
     \setlength{\ricol}{(\riinner-\ribridge)/2}%
     \begin{lrbox}{\ribox}\begin{minipage}{\riinner}%
        \centering
        {\color{darkblue}\scshape\large Research Interests}\\[4pt]
        \begin{minipage}{0.86\riinner}\centering\itshape #1\end{minipage}\\[9pt]
        \noindent\ignorespaces}
    {\end{minipage}\end{lrbox}%
     \noindent\makebox[\textwidth][r]{\vbox{\hsize=\rifull\offinterlineskip
        \hbox{\color{darkblue}\rule{\rifull}{1pt}}%
        \hbox{\colorbox{darkblue!6}{\usebox{\ribox}}}%
        \hbox{\color{darkblue}\rule{\rifull}{1pt}}}}%
     \par}

%%%%%%%%%%%%%%%%%%%%%%%% End Helper Commands %%%%%%%%%%%%%%%%%%%%%%%%%%%

%%%%%%%%%%%%%%%%%%%%%%%%% Begin CV Document %%%%%%%%%%%%%%%%%%%%%%%%%%%%

\begin{document}
\makeheading{Joshua Brendan Melander, PhD}

\section{Contact Information}

% NOTE: Mind where the & separators and \\ breaks are in the following
%       table. Table is one row made up of three parboxes. The left
%       parbox has address info, the middle parbox has a vertical bar,
%       and the right parbox has phone and electronic contact
%       information.
%
% MACROS: \rcollength is the width of the right column of the table
%             (adjust it to your liking; default is 1.85in).
%         \spacewidth is width of area between left and right boxes.
%
\newlength{\rcollength}\setlength{\rcollength}{1.85in}%
\newlength{\spacewidth}\setlength{\spacewidth}{20pt}
%
\begin{tabular}[t]{@{}p{\textwidth-\rcollength-\spacewidth}@{}p{\spacewidth}@{}p{\rcollength}@{}}%

% Address box
\parbox{\textwidth-\rcollength-\spacewidth}{%
	Pasca Lab\\
	Clark Center\\
	Stanford University, Stanford, CA 94305}
% \href{http://www.asu.edu/}{Arizona State University}\\
% \href{http://cidse.engineering.asu.edu/}{School of Computing and Augmented Intelligence}\\
% PO Box 878809, Room 553\\
% Tempe, AZ  85287-8809  USA}

&
% Uncomment to add a vertical bar in middle of contact information
% {\vrule width 0.5pt}
\parbox[m][5\baselineskip]{\spacewidth}{} &

% Non-snail-mail contact information
\parbox{\rcollength}{%
\textit{cell: } +1-669-666-4440 \\
\textit{email:} \email{melander@stanford.edu} \\
\textit{website: } \href{https://jbmelander.com}{jbmelander.com} \\
\textit{github: } \href{https://github.com/jbmelander}{jbmelander}} 


\end{tabular}

\begin{interests}{Trained in deep mechanistic modeling of biological circuits, I carry that rigor\\ across the boundary between living and artificial computation.}
\interest{Networks, read like neurons}{Biological interpretation of artificial neural networks}{Explaining what deep networks compute with the causal, circuit-level methods of neuroscience: sparse features and superposition, circuits and attribution, and new architectures extending and improving transformers.}%
\interestbridge%
\interest{Neurons, run like networks}{Biocomputing and neural modeling}{Engineering living human neural tissue into substrates that learn and compute: custom-wired circuits in neural organoids.}
\end{interests}

\section{Scientific Training}
\begin{lonelist}

	\item[] \textbf{Postdoc} Stanford University, California \hfill \textit{2026--present} 
        \begin{innerlist}
	\item \textit{PI}: Dr. Sergiu Pasca, Pasca Lab
	\item Computation $\times$ neural organoids
        \end{innerlist}

	\item[] \textbf{PhD} Stanford University, California \hfill \textit{2017--2025} 
        \begin{innerlist}
	\item Neuroscience Graduate Program
	\item \textit{Advisor}: Dr. Stephen Baccus
	\item Deep models of the retina and causal frameworks from neuroscience for AI interpretability
	\item \textit{Rotations}: Dr. Daniel Yamins, Dr. Karl Deisseroth, Dr. Stephen Baccus
        \end{innerlist}

	\item[] \textbf{Summer Course} Marine Biological Laboratory, Massachusetts \hfill \textit{2023}
        \begin{innerlist}
	\item Methods in Computational Neuroscience
	\item Studied recurrent-neural-network gating with Dr. Haim Sompolinsky
        \end{innerlist}

	\item[] \textbf{Post-Bac} Vollum Institute, Oregon \hfill \textit{2014--2017}
        \begin{innerlist}
	\item Research Assistant
	\item \textit{Advisors}: Dr. Haining Zhong and Dr. Tianyi Mao
        \end{innerlist}

\item[] \textbf{B.S.} Whitman College, Washington \hfill \textit{2010--2014}
        \begin{innerlist}
	\item Biochemistry, Biophysics, and Molecular Biology
	\item \textit{Magna Cum Laude}, Honors Thesis
	\item \textit{Advisors}: Dr. Christopher Wallace and Dr. Ginger Withers
        \end{innerlist}

\end{lonelist}


\section{Select Publications}
\begin{bibenum}
    \item[] \textbf{Melander, J.B.}, Alaoui, Z., Liu, S., Ganguli, S., Baccus, S. (2026) Causal interpretation of neural network computations with contribution decomposition (CODEC). \textit{International Conference on Learning Representations (ICLR)}. \href{https://jbmelander.com/melander2026_codec_iclr.pdf}{[pdf]} \href{https://github.com/baccuslab/bscope}{[code]}

    \item[] Ding, X., Lee, D., \textbf{Melander, J.B.}, Sivulka, G., Ganguli, S., Baccus, S. (2023) Information geometry of the retinal representation manifold. \textit{Advances in Neural Information Processing Systems (NeurIPS)}.

    \item[] Maheswaranathan, N., McIntosh, L., \ldots, \textbf{Melander, J.B.}, Nayebi, A., Brezovec, L., Wang, J., Ganguli, S., Baccus, S. (2023) Interpreting the retinal neural code for natural scenes: from computations to neurons. \textit{Neuron} 111, 2742--2755.

    \item[] \textbf{Melander, J.B.}*, Nayebi, A.*, Jongbloets, B., Fortin, D., Qin, M., Ganguli, S., Mao, T., Zhong, H. (2021) Distinct \textit{in vivo} dynamics of excitatory synapses onto cortical pyramidal neurons and parvalbumin-positive interneurons. \textit{Cell Reports} 37 (6), 109972. \href{https://jbmelander.com/melander2021_cellreports.pdf}{[pdf]}

    \item[] Maheswaranathan, N., McIntosh, L., Tanaka, H., Grant, S., Kastner, D., \textbf{Melander, J.B.}, Brezovec, L., Nayebi, A., Wang, J., Ganguli, S., Baccus, S. (2019) The dynamic neural code of the retina for natural scenes. \textit{bioRxiv} 430943.

    \item[] Fortin, D., \textbf{Melander, J.B.}, Jongbloets, B., Xiong, W., Guo, C., Mao, T., Zhong, H. (2019) High-contrast visualization of endogenous proteins for live imaging. \textit{Microscopy and Microanalysis} 25 (S2), 1250--1251.

    \item[] Ma, L., Jongbloets, B., Xiong, W., \textbf{Melander, J.B.}, Qin, M., Lameyer, T., Harrison, M., Zemelman, B., Mao, T., Zhong, H. (2018) A highly sensitive A-kinase activity reporter for imaging neuromodulatory events in awake mice. \textit{Neuron} 99 (4), 665--679.

    \item[] Fortin, D., Tillo, S., Yang, G., Rah, J., \textbf{Melander, J.B.}, Bai, S., Solar-Cedeno, O., Qin, M., Zemelman, B., Guo, C., Mao, T., Zhong, H. (2014) Live imaging of endogenous PSD-95 using ENABLED: a conditional strategy to fluorescently label endogenous proteins. \textit{Journal of Neuroscience} 34 (50), 16698--16712.
\end{bibenum}

\section{Open-Source Software}
\begin{lonelist}
\item[] \href{https://github.com/baccuslab/bscope}{\textbf{bscope}}: neural-network functional microscopy
	\begin{innerlist}
	\item PyTorch library for recording activations, gradients, and contributions in any model (CNNs, MLPs, ViTs), with attribution methods, causal ablations including attention-head disruption, and sparse autoencoders. Code for CODEC (ICLR 2026).
	\end{innerlist}
\end{lonelist}

\vspace{0.1in}

\section{Select Conference Presentations}
\begin{bibenum}
\item[] \textbf{Melander, J.B.}, Jolley, H., Miller, T., Withers, G., Wallace C. (2012) Coordinated cellular responses to experience-dependent plasticity. Murdock College Science Research Conference, Walla Walla, WA (Poster)
\item[] \textbf{Melander, J.B.}, Withers, G., Wallace C. (2013) Quantitative analysis of immunohistochemical markers of the induction of plasticity at the synapse and nucleus. Undergraduate poster session, Society for Neuroscience Annual Meeting, San Diego, CA (Poster)
\item[] \textbf{Melander, J.B.}, Jongbloets, B., Qin, M., Zhong, H., Mao, T. (2016) Imaging neuromodulation \textit{in vivo}. BRAIN Initiative Investigators Meeting, Bethesda, MD (Poster)
\item[] \textbf{Melander, J.B.}, Qin, M., Zhong, H., Mao, T. (2017) Chronic \textit{in vivo} imaging of excitatory shaft synapses onto cortical interneurons. Dendrite Gordon Research Conference, Lucca, Italy (Poster)
\item[] \textbf{Melander, J.B.}, Nayebi, A., Jongbloets, B., Fortin, D., Qin, M., Ganguli, S., Mao, T., Zhong, H. (2018) Chronic \textit{in vivo} imaging of excitatory synapses onto pyramidal and interneurons. Society for Neuroscience Conference, San Diego, CA (Poster)
\item[] Nayebi, A.*, \textbf{Melander, J.B.}*, Jongbloets, B., Mao, T., Zhong, H., Ganguli, S. (2019) Measuring and modeling the weight dynamics of many synapses onto diverse cell-types \textit{in vivo}. COSYNE, Lisbon, Portugal  (Talk)
\end{bibenum}

\section{Awards}
\begin{lonelist}
\item[ ] Alfred D. Abshire Scholar Research Award (2012)
\item[ ] Walter Brattain Academic Scholarship Award (2010--2014)
\item[ ] Oregon Developmental Biology Collaboration Award (2017)
\item[ ] Stanford Mind, Brain, Computation and Technology (Admittance, 2020)
\item[ ] Admitted to ``Methods in Computational Neuroscience'' summer session at Marine Biological Laboratory (2023)
\end{lonelist}

\section{Teaching}

\begin{lonelist}
\item[] Neurobiology (Biology 325), Whitman College (2013)
	\begin{innerlist}
	\item Laboratory TA: Organized staining, microscopy and electrophysiology exercises for students enrolled in Neurobiology labs
	\end{innerlist}
\item[] Genetics (Biology 303), Whitman College (2014)
	\begin{innerlist}
	\item Writing Fellow: Mentored other students taking Genetics in basic skills of scientific writing and figure preparation
	\end{innerlist}
\item[] Biophysics (BBMB 335), Whitman College (2014)
	\begin{innerlist}
	\item Laboratory TA: Organized computational exercises in protein folding simulation and visualization (Coot, QtMG), as well as wet-lab exercises including denaturation and biofilm assays
	\end{innerlist}
\item[] MATLAB and Microscopy Study Group, Vollum Institute (2016)
	\begin{innerlist}
	\item Lecturer: Taught general principles of classic machine-learning algorithms and implemented them in custom MATLAB scripts for a small study group of neurobiologists
	\end{innerlist}
\item[] CNJCx: Practical Python, Stanford University (2020)
	\begin{innerlist}
	\item Organizer, Director: A six-week curated virtual course designed to overcome the hurdles often faced by new programmers when approaching Python. Included numerous guest speakers, including professors and professional game designers. Supported by the Stanford Computational Neuroscience Journal Club and Stanford Mind, Brain, Computation, and Technology.
	\end{innerlist}




\end{lonelist}


\section{Mentorship}
\begin{lonelist}
\item[ ] Zaki Alaoui
    \begin{innerlist}
        \item \textit{Now}: REACH scholar at Stanford University
    \end{innerlist}
\item[ ] Prachi Shah
	\begin{innerlist}
	\item \textit{Then}: community-college intern in Baccus Lab
	\item \textit{Now}: undergraduate at UCI
	\end{innerlist}
\item[ ] Srinidhi Naidu
	\begin{innerlist}
	\item \textit{Then}: high-school student in Baccus Lab
	\item \textit{Now}: Undergraduate, University of Washington Neuroscience, Tuthill Lab
	\end{innerlist}
\end{lonelist}

\end{document}
